\chapter{Калибровочное усреднение цветных стрелок и межсекторный Краус-мост}
\label{ch:v8-gauge-twirl-cross-sector-kraus-bridge-gate}

\section{Почему линейный запрет не закрывает полную полугруппу}

Предыдущая глава выделила два центральных сектора
\begin{equation}
 \mathcal Z_{\rm fix}=\mathbb CP_q\oplus\mathbb CP_\ell,
 \qquad (\rank P_q,\rank P_\ell)=(12,9).
\end{equation}
Калибровочно-инвариантного линейного оператора между ними нет: кварковый сектор несёт
цветной триплет, а второй сектор является цветовым синглетом. Именно поэтому
предыдущий прямой соединитель вещественной структуры был уничтожен бесцветным проектором.

Но квантово-марковский генератор задаётся не одним инвариантным вектором, а
суммой по полному Краус-мультиплету. В физическом модуле стрелок уже имеются
две межсекторные семьи
\begin{equation}
 Q_L\longleftrightarrow Y_R,
 \qquad
 X_L\longleftrightarrow d_R,
 \label{eq:v8-existing-cross-arrow-multiplets}
\end{equation}
каждая размерности три над $\mathbb C$, или шесть над $\mathbb R$. Их
индивидуальные направления цветны, но вся двенадцатимерная вещественная
опора инвариантна как представление калибровочной группы.

\section{Базисно-независимая квадратичная сумма}

Пусть $\{e_a\}$ --- любой вещественный ортонормированный базис полной
межсекторной опоры стрелок, а
\begin{equation}
 D_{e_a}=\begin{pmatrix}0&e_a^*\\e_a&0\end{pmatrix}.
\end{equation}
Определим
\begin{equation}
 \mathcal L_{q\ell}
 =-\frac12\sum_a\operatorname{ad}_{D_{e_a}}^2.
 \label{eq:v8-gauge-twirled-kraus-generator}
\end{equation}
Сумма не зависит от выбора ортонормированного базиса. При калибровочном
преобразовании Краус-операторы смешиваются внутри полного мультиплета, а
генератор остаётся неизменным.

Машинный аудит проверил двенадцать случайных преобразований
$SU(3)\times SU(2)\times U(1)$; максимальная норма коммутатора
супероператоров равна $4.75\cdot10^{-15}$. Усреднение любой линейной
межсекторной стрелки по центру $SU(3)$ равно нулю с ошибкой
$1.66\cdot10^{-16}$. Следовательно, линейного цветового синглета и цветного
одноточечного среднего нет, а квадратичный процесс калибровочно-инвариантен.

\section{Снятие центральной двухсекторности}

На нормированном базисе $P_q/\sqrt{12},P_\ell/3$ ограничение положительной
формы $-\mathcal L_{q\ell}$ равно
\begin{equation}
 \begin{pmatrix}
 1&-2/\sqrt3\\
 -2/\sqrt3&4/3
 \end{pmatrix},
 \qquad
 \Spec=\left\{0,\frac73\right\}.
 \label{eq:v8-cross-kraus-central-restriction}
\end{equation}
По отдельности каждая из двух цветных семей уже имеет спектр
$\{0,7/6\}$ и сокращает $\mathbb C^2$ до скалярной единицы. Контрольные
семьи $L_L\leftrightarrow X_R$ и $Y_L\leftrightarrow e_R$, целиком лежащие
в бесцветном секторе, оставляют оба центральных направления нулевыми.

При независимом изменении положительных весов двух межсекторных семейств в
диапазоне $10^{-6}\ldots10^6$ все 64 теста дали одномерное центральное
ядро. Поэтому качественное снятие суперселекции не зависит от скорости
процесса, хотя сама скорость пока не выведена.

\begin{theorem}[Квадратичный межсекторный мост без цветного вакуума]
Полный физический модуль стрелок Тома VII содержит калибровочно-ковариантную
межсекторную Краус-опору. Её базисно-независимая квадратичная форма Дирихле
не имеет инвариантного линейного вектора, не требует цветного конденсата и
сокращает неподвижную алгебру
$\mathbb CP_q\oplus\mathbb CP_\ell$ до $\mathbb CI$. Следовательно,
двухсекторность не является строгой суперселекцией полного операторного
процесса.
\end{theorem}

\begin{remark}
Это не означает превращения отдельного кварка в лептон. Краус-сумма
описывает калибровочно ковариантный канал наблюдаемых, в котором заряд отдельного
оператора компенсируется полным мультиплетом. Физическая дилатация такого
канала и носитель компенсирующей информации ещё должны быть предъявлены.
\end{remark}

\begin{nogocriterion}[Запрет объявления полного родителя]
Нельзя отождествлять существование калибровочно-ковариантного Краус-моста с уже
выведенным взаимодействием наблюдаемых частиц. Следующий гейт обязан
встроить его в один кривизностный функционал, проверить гессианы Тома VII и
либо вывести относительную скорость, либо доказать независимость полного
качественного перехода от неё.
\end{nogocriterion}

\begin{protocol}
\begin{enumerate}
 \item межсекторных комплексных направлений стрелок: \textbf{$6$};
 \item вещественная Краус-опора: \textbf{$12$};
 \item линейный калибровочный синглет: \textbf{нет};
 \item квадратичная калибровочная ковариантность: \textbf{пройдена};
 \item цветное вакуумное среднее: \textbf{не требуется};
 \item центральное ядро после моста: \textbf{$\mathbb CI$};
 \item устойчивость по положительным скоростям: \textbf{пройдена};
 \item единственная скорость: \textbf{не выведена};
 \item общий родительского действия и гессиан: \textbf{открыты};
 \item строгая суперселекция $\mathbb C^2$: \textbf{опровергнута}.
\end{enumerate}
\end{protocol}

Машинный сертификат сохранён в
\path{s2t_v8_gauge_twirl_cross_sector_kraus_bridge_gate_results.json}.

\begin{thebibliography}{9}
\bibitem{v8-kraus-vacchini}
B.~Vacchini,
\emph{Covariant Mappings for the Description of Measurement, Dissipation
and Decoherence in Quantum Mechanics}, arXiv:0707.0603.
\bibitem{v8-kraus-ritter}
W.~G.~Ritter,
\emph{Quantum Channels and Representation Theory},
arXiv:quant-ph/0502153.
\bibitem{v8-kraus-kuramochi}
Y.~Kuramochi,
\emph{Nonstandard derivation of the Gorini--Kossakowski--Sudarshan--Lindblad
master equation from the Kraus representation}, arXiv:2406.03775.
\end{thebibliography}